\documentclass[12pt,a4paper]{article}
\usepackage[utf8]{inputenc}
\usepackage{amsmath,amssymb,amsthm}
\usepackage{geometry}
\usepackage{hyperref}
\usepackage{algorithm}
\usepackage{algpseudocode}
\usepackage{booktabs}

\geometry{margin=2.5cm}
\newtheorem{theorem}{Theorem}
\newtheorem{lemma}[theorem]{Lemma}
\newtheorem{corollary}[theorem]{Corollary}
\newtheorem{definition}{Definition}

\title{Problem 453: Lattice Quadrilaterals}
\author{}
\date{}

\begin{document}
\maketitle

\section{Problem Statement}

A \emph{simple quadrilateral} is a polygon with four distinct vertices, no straight angles, and no self-intersections.

Let $Q(m, n)$ be the number of simple quadrilaterals whose vertices are lattice points $(x, y)$ with $0 \le x \le m$ and $0 \le y \le n$.

Given: $Q(3, 7) = 39590$, $Q(12, 3) = 309000$, $Q(123, 45) = 70542215894646$.

\textbf{Find} $Q(12345, 6789) \bmod 135707531$.

\section{Mathematical Foundation}

Let $N = (m+1)(n+1)$ denote the total number of lattice points in the grid.

\begin{lemma}[Convex case]\label{lem:convex}
For any 4 points in convex position with no 3 collinear, exactly 1 of the 3 possible cyclic orderings yields a simple quadrilateral.
\end{lemma}

\begin{proof}
Let $A, B, C, D$ be in convex position with hull order $ABCD$. The three pairings of opposite edges give orderings $ABCD$, $ABDC$, $ACBD$. For $ABCD$, edges follow the convex hull, so no crossing occurs. For the other two orderings, a pair of ``opposite'' edges consists of diagonals of the convex hull, which must intersect by the convexity of the quadrilateral. Hence exactly one ordering is simple.
\end{proof}

\begin{lemma}[Concave case]\label{lem:concave}
For any 4 points in concave position (one strictly inside the triangle of the other 3), with no 3 collinear, all 3 cyclic orderings yield simple quadrilaterals.
\end{lemma}

\begin{proof}
Let $D$ lie strictly inside $\triangle ABC$. In any cyclic ordering, edges incident to $D$ are contained in the interior of $\triangle ABC$, while the edge connecting two of $\{A, B, C\}$ lies on or outside the triangle. A case analysis of the three orderings confirms no pair of non-adjacent edges can intersect.
\end{proof}

\begin{lemma}[No double-counting]\label{lem:nocross}
If a 4-point set contains two distinct collinear triples, then all 4 points are collinear.
\end{lemma}

\begin{proof}
Two 3-element subsets of a 4-element set share exactly 2 points, which determine a unique line. Both collinear triples lie on this line, so all 4 points are collinear.
\end{proof}

\begin{theorem}[Main formula]\label{thm:main}
\[
Q(m,n) = \binom{N}{4} - D + 2 \sum_T I(T),
\]
where $D$ is the number of degenerate 4-subsets, and $\sum_T I(T)$ sums the interior lattice points over all non-degenerate lattice triangles~$T$.
\end{theorem}

\begin{proof}
Let $S = \binom{N}{4} - D$ be the count of non-degenerate 4-subsets. Partition into convex ($S_{\text{cvx}}$) and concave ($S_{\text{ccv}}$). By Lemmas~\ref{lem:convex} and~\ref{lem:concave}:
\[
Q = S_{\text{cvx}} + 3 S_{\text{ccv}} = S + 2 S_{\text{ccv}}.
\]
Each concave 4-subset $\{A,B,C,D\}$ with $D$ inside $\triangle ABC$ corresponds uniquely to the pair $(\triangle ABC, D)$, so $S_{\text{ccv}} = \sum_T I(T)$. Thus $Q = \binom{N}{4} - D + 2\sum_T I(T)$.
\end{proof}

\begin{theorem}[Degenerate count]\label{thm:degen}
\[
D = \sum_{L} \left[\binom{k_L}{4} + \binom{k_L}{3}(N - k_L)\right],
\]
where the sum runs over all lines $L$ with $k_L \ge 3$ grid points.
\end{theorem}

\begin{proof}
$\binom{k_L}{4}$ counts 4-subsets entirely on~$L$; $\binom{k_L}{3}(N-k_L)$ counts those with exactly 3 on~$L$ and 1 off. By Lemma~\ref{lem:nocross}, no double-counting occurs.
\end{proof}

\begin{theorem}[Interior sum via Pick's theorem]\label{thm:pick}
By Pick's theorem, $I(T) = A(T) - B(T)/2 + 1$. Summing over non-degenerate triangles:
\[
\sum_T I(T) = \sum_T A(T) - \frac{1}{2}\sum_T B(T) + T_{\mathrm{count}},
\]
where $T_{\mathrm{count}} = \binom{N}{3} - \sum_L \binom{k_L}{3}$.
\end{theorem}

\begin{proof}
Sum Pick's theorem $I(T) = A(T) - B(T)/2 + 1$ over all non-degenerate triangles; the constant~$1$ sums to $T_{\mathrm{count}}$.
\end{proof}

\section{Algorithm}

\begin{algorithm}[H]
\caption{Compute $Q(m, n) \bmod p$}
\begin{algorithmic}[1]
\State Enumerate all primitive directions $(a, b)$ with $\gcd(a, |b|) = 1$, $a \ge 0$.
\For{each direction $(a, b)$}
    \State Compute line counts $k_L$ for each parallel line through the grid.
    \State Accumulate $D$: $\sum \binom{k_L}{4} + \binom{k_L}{3}(N - k_L)$.
    \State Accumulate $\sum_T B(T)$ via GCD sums: for a line with $k$ points, the pair-GCD sum is $k(k^2-1)/6$.
    \State Accumulate $\sum_T A(T)$ via signed-distance sums $F(c) = \sum_{(x,y)} |bx - ay - c|$, computed with prefix sums.
\EndFor
\State Combine: $Q = \binom{N}{4} - D + 2\bigl[\sum_T A(T) - \tfrac{1}{2}\sum_T B(T) + T_{\mathrm{count}}\bigr] \pmod{p}$
\State \Return $Q$
\end{algorithmic}
\end{algorithm}

\section{Complexity Analysis}

\begin{itemize}
    \item \textbf{Time:} $O(mn(m+n))$. There are $O(mn)$ primitive directions, each requiring $O(m+n)$ work for line statistics and prefix-sum area computation.
    \item \textbf{Space:} $O(m+n)$ for prefix sums per direction.
\end{itemize}

\section{Answer}

\[
\boxed{104354107}
\]

\end{document}
